\subsection*{Note (2)}
The second note, beginning on p.~52 and ending on p.~53, qualifies the
claim in \S3, p.~12, that no mathematical fact had made an impression
comparable to the discovery attached to a child's drawing. Grothendieck
recalls learning the definition of a circle from Maria, an internee
who gave him lessons at Rieucros when he was about twelve. The definition's
simplicity and obviousness transformed what had seemed the mysterious
``rotondité parfaite'' of the circle. He identifies this retrospectively
as his first glimpse of the creative power of a good definition, a
``formulation qui décrit l'essence'', and says that its fascination has
persisted. The note thus supplies the explicit earlier exception to the
turning point narrated in \S3 while returning to the role of finding a
formulation that the opening discussion of writing in \S1 had already
emphasized.
(\href{https://grothendieck.umontpellier.fr/web/viewer.html?file=../119.pdf\#page=55}{Copy 1, pp.~52--53, beginning at site 54};
\href{https://grothendieck.umontpellier.fr/web/viewer.html?file=../119.pdf\#page=56}{ending at site 55};
\href{https://grothendieck.umontpellier.fr/web/viewer.html?file=../119.pdf\#page=15}{call in \S3, p.~12, site 14}.)
